% Generated by scripts/assemble.py; corrections belong in source/corrections.json.
\examhead{14}{北京邮电大学 2024 年数学分析试卷}{19}
% p020-b06; source PDF page 20
\par\Needspace{5\baselineskip}
\par\Needspace{6\baselineskip}
\blockanchor{p020-b06}
\majorquestion{一}求极限\par
\nopagebreak[4]
% p020-b07; source PDF page 20
\blockanchor{p020-b07}
\displayformula{\lim _ {x \to 0} \frac {1 - \cos x \sqrt {\cos 2 x} \sqrt [ 3 ]{\cos 3 x}}{x ^ {2}}.}
% p020-b08; source PDF page 20
\par\Needspace{5\baselineskip}
\par\Needspace{6\baselineskip}
\blockanchor{p020-b08}
\majorquestion{二}计算\par
\nopagebreak[4]
% p020-b09; source PDF page 20
\blockanchor{p020-b09}
\displayformula{\int_ {- \frac {\pi}{2}} ^ {\frac {\pi}{2}} \frac {1 + \sin 3 x}{1 + \cos^ {2} x} \mathrm{d} x.}
% p020-b10; source PDF page 20
\par\Needspace{5\baselineskip}
\blockanchor{p020-b10}
\majorquestion{三}求幂级数 $\sum_{n=1}^{\infty} \frac{(-1)^{(n-1)}}{2n-1} x^{2n}$ 的收敛域及和函数.\par
% p020-b11; source PDF page 20
\par\Needspace{5\baselineskip}
\blockanchor{p020-b11}
\majorquestion{四}求二元函数 $f(x, y) = x^{2}(2 + y^{2}) + y \ln y$ 的极值.\par
% p020-b12; source PDF page 20
\par\Needspace{5\baselineskip}
\blockanchor{p020-b12}
\majorquestion{五}若 $a_{n} > 0$, $\lim_{n\to \infty}\frac{a_n}{a_{n + 1}} = l > 1$, 证明: $\lim_{n\to \infty}a_n = 0$.\par
% p020-b13; source PDF page 20
\par\Needspace{5\baselineskip}
\blockanchor{p020-b13}
\majorquestion{六}讨论 $\ln^2 x - 2\ln x + 2x + k = 0$ 的实根个数, 其中 $k$ 是常数.\par
% p020-b14; source PDF page 20
\par\Needspace{5\baselineskip}
\blockanchor{p020-b14}
\majorquestion{七}已知 $x \mathrm{e}^{-\frac{x}{2}} - \ln (1 + x)$ 与 $\alpha x^n$ 是等价无穷小量, 求 $\alpha$ 和 $n$.\par
% p020-b15; source PDF page 20
\par\Needspace{5\baselineskip}
\blockanchor{p020-b15}
\majorquestion{八}判别 $\int_0^{+\infty} \frac{\ln x}{\sqrt[3]{x^4 - x}} \, \mathrm{d}x$ 的敛散性.\par
% p020-b16; source PDF page 20
\par\Needspace{5\baselineskip}
\blockanchor{p020-b16}
\majorquestion{九}求圆锥体 $z \geq \sqrt{x^2 + y^2}$ 与 $x^2 + y^2 + (z - 1)^2 \leq 1$ 所围立体的体积.\par
% p020-b17; source PDF page 20
\par\Needspace{5\baselineskip}
\par\Needspace{6\baselineskip}
\blockanchor{p020-b17}
\majorquestion{十}计算\par
\nopagebreak[4]
% p020-b18; source PDF page 20
\blockanchor{p020-b18}
\displayformula{I = \int_ {L} 3 x y \mathrm{d} x + (x ^ {3} + x - 2 y) \mathrm{d} y.}
% p020-b19; source PDF page 20
\blockanchor{p020-b19}
其中 $L$ 是第一象限从 $(0,0)$ 沿圆周 $x^{2} + y^{2} - 2x = 0$ 到 $(2,0)$，再从 $(2,0)$ 沿圆周 $x^{2} + y^{2} = 4$ 到 $(0,2)$ 的曲线.\par
% p020-b20; source PDF page 20
\par\Needspace{5\baselineskip}
\par\Needspace{6\baselineskip}
\blockanchor{p020-b20}
\majorquestion{十一}平面 $2x + y + 2z = 2$ 与三个坐标平面围成 $\Omega, \Sigma$ 为 $\Omega$ 的整个表面, 取外侧, 求\par
\nopagebreak[4]
% p020-b21; source PDF page 20
\blockanchor{p020-b21}
\displayformula{I = \iint_ {\Sigma} (x ^ {2} + 1) \mathrm{d} y \mathrm{d} z - 2 y \mathrm{d} z \mathrm{d} x + 3 z \mathrm{d} x \mathrm{d} y.}
% p020-b22; source PDF page 20
\par\Needspace{5\baselineskip}
\blockanchor{p020-b22}
\majorquestion{十二}讨论 $\alpha > 0$ 时, $f(x) = x^{\alpha}$ 在 $(0, +\infty)$ 上的一致连续性.\par
% p020-b23; source PDF page 20
\par\Needspace{5\baselineskip}
\blockanchor{p020-b23}
\majorquestion{十三}证明 $\sum_{n=1}^{\infty} \frac{(-1)^n}{x^2 + n}$ 在 $(- \infty, +\infty)$ 上一致收敛.\par
% p020-b24; source PDF page 20
\par\Needspace{5\baselineskip}
\blockanchor{p020-b24}
\majorquestion{十四}$f(x)$ 在 $(0, +\infty)$ 上可微, $0 \leq f(x) \leq \frac{x}{1 + x^2}$, 证明存在 $\xi > 0$, 使 $f'(\xi) = \frac{1 - \xi^2}{(1 + \xi^2)^2}$.\par
% p020-b25; source PDF page 20
\par\Needspace{5\baselineskip}
\blockanchor{p020-b25}
\majorquestion{十五}设  $f(x)$  可导且  $f(0)=1, 0 < f'(x) < \frac{1}{2}$ ，数列  $\{x_{n}\}$  满足  $x_{n+1}=f(x_{n}) (n=1,2,\cdots)$ . 证明:\par
% p020-b26; source PDF page 20
\blockanchor{p020-b26}
\question{1}级数 $\sum_{n=1}^{\infty}(x_{n+1} - x_n)$ 绝对收敛.\par
% p020-b27; source PDF page 20
\blockanchor{p020-b27}
\question{2}$\lim_{n\to \infty}x_n$ 存在且 $0 < \lim_{n\to \infty}x_n < 2.$\par
